% MIXED: actual event labels and proved asymptotic scales; schematic spacing.
% Arrows mean actual event passages (top) or discrete returns (bottom).
\begin{figure}[!htbp]
\centering
\begin{tikzpicture}[
 x=1cm,y=1cm,font=\small,
 every node/.style={inner sep=2pt},
 stage/.style={draw=black!65,rounded corners=2pt,align=center,
               text width=2.35cm,minimum height=1.40cm},
 arrow/.style={-{Latex[length=1.6mm,width=1.1mm]},line width=.5pt},
 chain/.style={draw=black!65,align=center,text width=4.3cm,
               minimum height=.75cm},
 guide/.style={draw=black!40,densely dotted,line width=.5pt}
]
 \node[anchor=west,font=\small\bfseries] at (0,7.05)
   {(a) Actual events from the physical initial state};
 \node[stage] (initial) at (1.25,5.45)
   {original\\preparation};
 \node[stage] (release) at (4.25,5.45)
   {critical release\\$v_r/\varepsilon\to q_{\rm in}>0$};
 \node[stage] (tdone) at (7.25,5.45)
   {signed $\mathrm{TD}_1$\\and contact\\$p_1/\sqrt\varepsilon\to\pi$};
 \node[stage] (tdtwo) at (10.25,5.45)
   {complete $\mathrm{TD}_2$\\$\delta_2\sim C_R\varepsilon^{3/8}$};
 \node[stage] (tdthree) at (13.25,5.45)
   {admitted $\mathrm{TD}_3$\\$\delta_3\asymp\varepsilon^{9/32}$};
 \draw[arrow] (initial.east) -- (release.west);
 \draw[arrow] (release.east) -- (tdone.west);
 \draw[arrow] (tdone.east) -- (tdtwo.west);
 \draw[arrow] (tdtwo.east) -- (tdthree.west);
 \node at (5.75,6.55) {first flight};
 \node at (8.75,6.55) {second flight};
 \node at (11.75,6.55) {one return};
 \node[anchor=west] at (0,4.40) {$\pi$ may have either sign or vanish.};

 \node[anchor=west,font=\small\bfseries] at (0,3.65)
   {(b) A comparison started at the same actual base};
 \node[chain] (actualstart) at (2.35,2.82)
   {actual $(\delta_3,\xi_3^r)$};
 \node[chain] (graphstart) at (2.35,1.65)
   {graph $(\delta_3,\Gamma_\lambda(\delta_3))$};
 \node[chain] (actualend) at (12.05,2.82)
   {actual $(\delta_j,\xi_j^r)$};
 \node[chain] (graphend) at (12.05,1.65)
   {graph $(\bar\delta_j,\Gamma_\lambda(\bar\delta_j))$};
 \draw[arrow] (actualstart.east) -- (6.15,2.82);
 \node at (7.20,2.82) {$\cdots$};
 \draw[arrow] (8.25,2.82) -- (actualend.west);
 \draw[arrow,densely dashed] (graphstart.east) -- (6.15,1.65);
 \node at (7.20,1.65) {$\cdots$};
 \draw[arrow,densely dashed] (8.25,1.65) -- (graphend.west);
 \draw[guide] (actualstart.south) -- (graphstart.north);
 \node[fill=white,align=center] at (7.20,2.23)
   {same event indices;\\unbounded return count};
 \node[anchor=west] at (0,.82)
   {Same actual base; different complete fibres.};
 \node[anchor=west] at (0,.22)
   {Same-index comparison: Lemma~\ref{sel:sharp-phase}.};
\end{tikzpicture}
\caption{From the physical preparation to the same-base return
comparison. (a) Both preparatory flights and the intervening signed
contact are generated by the original trajectory.
(b) The graph starts at the actual third-touchdown base, with a small
complete fibre discrepancy. The actual normal memory is retained;
the two starting states are not identical.
Arrows and spacing are schematic, not physical time.}
\label{fig:entry-comparison}
\end{figure}
